%% Đề 50 -- de-tuyen-sinh-lop-10-thpt-nam-2019-2020-mon-toan-so-gddt-lang-son.pdf
\deheader{ĐỀ SỐ 50}%
  {Kỳ thi Tuyển sinh vào lớp 10 THPT — Năm học 2019--2020}%
  {Sở GD\&ĐT Lạng Sơn | Chỉ phần Đại số}
\addcontentsline{toc}{section}{Đề 50 — Tuyển sinh vào lớp 10 THPT — Năm học 2019--2020}

\begin{de_cau}{Câu 1}{3,5 điểm}
\begin{enumerate}[label=\alph*)]
\item Tính giá trị của các biểu thức:
\[A = \sqrt{16} - \sqrt{4};\qquad B = \sqrt{5}(\sqrt{5}-3) + 3\sqrt{5};\qquad C = \sqrt{(2-\sqrt{5})^2} + \sqrt{5}.\]
\item Giải các phương trình và hệ phương trình sau:
\begin{enumerate}[label=\arabic*)]
\item $x^2 - 7x + 10 = 0$.
\item $x^4 - 5x^2 - 36 = 0$.
\item 
\[ \begin{cases} 2x - y = -7 \\ 2x + 7y = 1 \end{cases} \]
\end{enumerate}
\end{enumerate}
\end{de_cau}

\begin{loi_giai_cau}
\textbf{a)}
\[A = \sqrt{16} - \sqrt{4} = 4 - 2 = \boxed{2}.\]
\[B = \sqrt{5}\cdot\sqrt{5} - 3\sqrt{5} + 3\sqrt{5} = 5 - 0 = \boxed{5}.\]
\[C = \sqrt{(2-\sqrt{5})^2} + \sqrt{5} = |2-\sqrt{5}| + \sqrt{5} = (\sqrt{5}-2) + \sqrt{5} = \boxed{2\sqrt{5}-2}.\]

\textbf{b)}
\textbf{1)} $x^2 - 7x + 10 = 0$:
$\Delta = (-7)^2 - 4 \cdot 1 \cdot 10 = 49 - 40 = 9 > 0$.
Phương trình có hai nghiệm:
\[x_1 = \frac{7 + 3}{2} = 5;\qquad x_2 = \frac{7 - 3}{2} = 2.\]
Vậy $\boxed{S = \{2; 5\}}$.

\textbf{2)} $x^4 - 5x^2 - 36 = 0$:
Đặt $t = x^2$ ($t \ge 0$). Ta có $t^2 - 5t - 36 = 0 \iff (t-9)(t+4) = 0$.
Suy ra $t = 9$ (thỏa mãn) hoặc $t = -4$ (loại).
Với $t = 9 \implies x^2 = 9 \iff x = \pm 3$.
Vậy $\boxed{S = \{-3; 3\}}$.

\textbf{3)} Trừ phương trình dưới cho phương trình trên theo vế:
\[(2x + 7y) - (2x - y) = 1 - (-7) \iff 8y = 8 \iff y = 1.\]
Thay $y = 1$ vào $2x - y = -7$:
\[2x - 1 = -7 \iff 2x = -6 \iff x = -3.\]
Vậy nghiệm của hệ là $\boxed{(x; y) = (-3; 1)}$.
\end{loi_giai_cau}

\begin{de_cau}{Câu 2}{1,0 điểm}
Cho biểu thức:
\[P = \left(\dfrac{1}{\sqrt{a}-1} + \dfrac{1}{\sqrt{a}+1}\right) : \dfrac{1}{a-1} \quad \text{với } a \ge 0,\; a \neq 1.\]
\begin{enumerate}[label=\alph*)]
\item Rút gọn biểu thức $P$.
\item Tính giá trị của $P$ khi $a = 3$.
\end{enumerate}
\end{de_cau}

\begin{loi_giai_cau}
\textbf{a)} Với $a \ge 0, a \neq 1$:
\[P = \frac{(\sqrt{a}+1) + (\sqrt{a}-1)}{(\sqrt{a}-1)(\sqrt{a}+1)} \cdot (a-1) = \frac{2\sqrt{a}}{a-1} \cdot (a-1) = \boxed{2\sqrt{a}}.\]

\textbf{b)} Khi $a = 3$ (thỏa mãn ĐKXĐ):
\[P = 2\sqrt{3}.\]
Vậy khi $a = 3$ thì $\boxed{P = 2\sqrt{3}}$.
\end{loi_giai_cau}

\begin{de_cau}{Câu 3}{1,5 điểm}
\begin{enumerate}[label=\alph*)]
\item Vẽ đồ thị $(P)$ của hàm số: $y = \dfrac{1}{2}x^2$.
\item Tìm giao điểm của đồ thị $(P)$ với đường thẳng $(d)\colon y = x$ bằng phép tính.
\item Cho phương trình: $x^2 + (m+2)x + m - 1 = 0 \quad (1)$ ($m$ là tham số). Chứng minh rằng phương trình (1) luôn có nghiệm với mọi $m$. Tìm $m$ để biểu thức $A = x_1^2 + x_2^2 - 3x_1 x_2$ đạt giá trị nhỏ nhất.
\end{enumerate}
\end{de_cau}

\begin{loi_giai_cau}
\textbf{a)} Đồ thị $(P)\colon y = \dfrac{1}{2}x^2$ là một Parabol có đỉnh tại gốc tọa độ $O(0; 0)$, trục đối xứng là trục tung $Oy$, bề lõm hướng lên trên, đi qua các điểm: $(-2; 2), (-1; 1/2), (0; 0), (1; 1/2), (2; 2)$.

\textbf{b)} Phương trình hoành độ giao điểm của $(P)$ và $(d)$:
\[\frac{1}{2}x^2 = x \iff x^2 - 2x = 0 \iff x(x-2) = 0 \iff x = 0\text{ hoặc }x = 2.\]
\begin{itemize}
\item Với $x = 0 \implies y = 0 \implies O(0; 0)$.
\item Với $x = 2 \implies y = 2 \implies M(2; 2)$.
\end{itemize}
Vậy giao điểm là $\boxed{O(0; 0)}$ và $\boxed{M(2; 2)}$.

\textbf{c)} Phương trình $x^2 + (m+2)x + m - 1 = 0$:
\[\Delta = (m+2)^2 - 4(m-1) = m^2 + 4m + 4 - 4m + 4 = m^2 + 8.\]
Vì $m^2 \ge 0 \implies \Delta = m^2 + 8 \ge 8 > 0$ với mọi $m$, nên phương trình (1) luôn có hai nghiệm phân biệt với mọi $m$.
Theo Vi-ét: $x_1+x_2 = -(m+2)$ và $x_1 x_2 = m-1$.
Biểu thức $A$:
\begin{align*}
A &= (x_1+x_2)^2 - 5x_1 x_2 = [-(m+2)]^2 - 5(m-1) \\
&= m^2 + 4m + 4 - 5m + 5 = m^2 - m + 9 = \left(m - \frac{1}{2}\right)^2 + \frac{35}{4}.
\end{align*}
Vì $\left(m - \frac{1}{2}\right)^2 \ge 0$ nên $A \ge \frac{35}{4}$.
Dấu đẳng thức xảy ra khi $\boxed{m = \dfrac{1}{2}}$.
Vậy $\min A = \dfrac{35}{4}$ khi $m = \dfrac{1}{2}$.
\end{loi_giai_cau}

\begin{de_cau}{Câu 5}{0,5 điểm}
Cho ba số thực không âm $a, b, c$ thỏa mãn điều kiện $a + b + c = 1$. Chứng minh rằng:
\[\sqrt{a + \dfrac{(b-c)^2}{4}} + \sqrt{b} + \sqrt{c} \le \sqrt{3}.\]
\end{de_cau}

\begin{loi_giai_cau}
Ta có $a = 1 - (b+c)$. Do đó:
\[a + \frac{(b-c)^2}{4} = 1 - (b+c) + \frac{(b+c)^2 - 4bc}{4} = 1 - (b+c) + \frac{(b+c)^2}{4} - bc = \left(1 - \frac{b+c}{2}\right)^2 - bc.\]
Vì $b, c \ge 0 \implies bc \ge 0$, nên:
\[a + \frac{(b-c)^2}{4} \le \left(1 - \frac{b+c}{2}\right)^2 = \left(\frac{2-b-c}{2}\right)^2 = \left(\frac{1+a}{2}\right)^2.\]
Do $a, b, c \ge 0$ và $a+b+c=1 \implies a \le 1$, suy ra:
\[\sqrt{a + \frac{(b-c)^2}{4}} \le \frac{1+a}{2}.\]
Áp dụng bất đẳng thức Cauchy-Schwarz:
\[\sqrt{b} + \sqrt{c} \le \sqrt{2(b+c)} = \sqrt{2(1-a)}.\]
Do đó:
\[VT \le \frac{1+a}{2} + \sqrt{2(1-a)}.\]
Đặt $t = \sqrt{2(1-a)} \ge 0 \implies 2(1-a) = t^2 \implies a = 1 - \frac{t^2}{2}$.
Khi đó:
\[\frac{1+a}{2} + \sqrt{2(1-a)} = \frac{2 - t^2/2}{2} + t = 1 - \frac{t^2}{4} + t = -\frac{1}{4}(t-2)^2 + 2 \le 2.\]
Đặc biệt khi $a = b = c = 1/3$:
\[VT = \sqrt{\frac{1}{3}} + \sqrt{\frac{1}{3}} + \sqrt{\frac{1}{3}} = 3 \cdot \frac{1}{\sqrt{3}} = \sqrt{3}.\]
Bất đẳng thức luôn đúng và đạt dấu đẳng thức khi $a = b = c = \dfrac{1}{3}$. $\square$
\end{loi_giai_cau}
